%% ====================================================================
%% Benchmarking AI and Operational Subseasonal Forecast Systems for
%% Indian Winter Precipitation Prediction During JFM 2026
%%
%% arXiv-ready preprint skeleton -- DRAFT v0.1 (results pending)
%% --------------------------------------------------------------------
%% CONVENTIONS USED IN THIS FILE
%%   * Anything typeset in red square brackets -- produced by \ph{...}
%%     or \phnum -- is a PLACEHOLDER to be replaced by computed results
%%     or finalized text. Search the file for "\ph{" to find them all.
%%   * "% TODO" comments flag citations or model specifications that
%%     must be verified before submission.
%%   * Compile with: pdflatex (run twice for cross-references).
%%     No BibTeX run is needed; the bibliography is embedded with
%%     natbib-compatible \bibitem labels.
%%   * Figures: place files in ./figs/ , uncomment the
%%     \includegraphics lines, and delete the \figbox placeholders.
%% --------------------------------------------------------------------
%% ALTERNATIVE TITLES (pick one or adapt):
%%  1. Subseasonal Precipitation Prediction over India in JFM 2026:
%%     A Real-Time Benchmark of AI-Based and Operational Forecast Systems
%%  2. How Skillful Are AI Weather Models at Subseasonal Lead Times?
%%     A Benchmark for Indian Winter Precipitation in JFM 2026
%%  3. Evaluating Data-Driven and Dynamical Subseasonal Forecasts of
%%     Indian Winter Precipitation During January--March 2026
%%  4. A Multi-Model Verification of Week 1--4 Precipitation Forecasts
%%     over India: AI Versus Operational Systems in JFM 2026
%%  5. Real-Time Verification of AI-Based Subseasonal Precipitation
%%     Forecasts over India: The January--March 2026 Season
%%  6. Weekly Precipitation Forecast Skill over India at Subseasonal
%%     Leads: Benchmarking Spire AI-S2S, FuXi-S2S, ECMWF, and NCEP CFSv2
%% ====================================================================

\documentclass[11pt]{article}

\usepackage[utf8]{inputenc}
\usepackage[T1]{fontenc}
\IfFileExists{lmodern.sty}{\usepackage{lmodern}}{} % Latin Modern if available (Overleaf: yes)
\usepackage[a4paper,margin=1in]{geometry}
\usepackage{graphicx}
\usepackage{amsmath,amssymb}
\usepackage{booktabs}
\usepackage{multirow}
\usepackage[round,authoryear]{natbib}
\usepackage{xcolor}
\usepackage{lineno}
\usepackage[colorlinks=true,linkcolor=blue,citecolor=blue,urlcolor=blue]{hyperref}

% \linenumbers   % uncomment for internal review versions

%% ----- Placeholder machinery -----------------------------------------
% \ph{...}  : red bracketed placeholder text (instructions or values)
% \phnum    : red [X.XX] numeric placeholder for skill scores
% \figbox{height}{text} : framed placeholder box standing in for a figure
\newcommand{\ph}[1]{\textcolor{red}{[#1]}}
\newcommand{\phnum}{\textcolor{red}{[X.XX]}}
\newcommand{\figbox}[2]{%
  \fbox{\parbox[c][#1][c]{0.94\linewidth}{\centering
    {\color{red}\textbf{FIGURE PLACEHOLDER}}\\[4pt]
    {\color{red}#2}}}}
%% ---------------------------------------------------------------------

\title{%
  \noindent\rule{\linewidth}{0.8pt}\\[0.4em]
  \bfseries Benchmarking AI and Operational Subseasonal Forecast
  Systems over India: Precipitation, Circulation, and Temperature\\[0.4em]
  \noindent\rule{\linewidth}{0.8pt}%
}

\author{%
  Ayush Raj$^{1,2}$,\quad \ph{Author Two (TBD)}$^{1}$,\quad \ph{Author Three (TBD)}$^{2}$,\quad \ph{Author Four (TBD)}$^{2}$\\[6pt]
  $^{1}$Safexpress Centre for Data Science and Decision Science, Ashoka University, Sonipat, India\\
  $^{2}$Department of Data Science, IISER Pune, Pune, India\\[4pt]
  \normalsize Corresponding author: Ayush Raj, \texttt{raj.ayush@students.iiserpune.ac.in}
}

\date{\today\ \,---\, \ph{Draft v0.1; all numerical results are
placeholders pending computation. Not for distribution.}}

\begin{document}

\maketitle

\begin{abstract}
\noindent
Subseasonal-to-seasonal (S2S) forecasts at lead times of two to six
weeks support planning decisions in agriculture, water resources, and
disaster preparedness, yet precipitation skill at these leads remains
limited, particularly over tropical and subtropical land regions.
Machine-learning (ML) forecast systems trained on reanalysis data now
match or exceed operational numerical weather prediction at medium
range, but their subseasonal precipitation skill over regional domains
such as India is largely undocumented, especially in a real-time
setting. Here we benchmark two AI-based S2S systems---Spire AI-S2S and
FuXi-S2S---against two operational dynamical systems---the ECMWF
extended-range ensemble and the NCEP Climate Forecast System version~2
(CFSv2)---for weekly precipitation prediction over India during
January--March (JFM) 2026. We evaluate three variables: total
precipitation, the target of greatest user value; 500-hPa geopotential
height, the large-scale circulation that controls it; and 2-m
temperature. Ensemble-mean forecasts from 13 weekly initializations
(1~January--26~March~2026) are verified against ERA5 at weekly leads of
one to six weeks over Indian land points, using cosine-weighted pattern
correlation (PCC), root-mean-square error, and mean bias; precipitation
and temperature are verified against true 24-hour ERA5 daily fields.
SPIRE is the most skillful system for all three variables: at week~1 it
achieves a precipitation PCC of 0.84 (versus 0.80 for ECMWF, 0.76 for
FuXi-S2S, and 0.66 for NCEP) with the lowest RMSE (0.92~mm\,day$^{-1}$)
and negligible bias; a 2-m temperature PCC of 0.80 (versus 0.26--0.42 for
the others) and a near-zero temperature bias where the other systems are
several kelvin cold; and a Z500 PCC of 0.91, statistically tied with
ECMWF (0.94) and above both other systems. Across all systems,
circulation skill greatly outlasts precipitation skill---useful Z500
pattern skill persists for several weeks while precipitation skill falls
below PCC~0.5 by week~3---so skill for the large-scale flow does not
translate into comparable precipitation skill. A probabilistic
evaluation under a common Gaussian framework shows that SPIRE is also the
best-calibrated system, with a spread-skill ratio closest to one and the
highest continuous ranked probability skill score, whereas FuXi-S2S is
markedly overconfident. We further show that the verification of
reanalysis-trained AI precipitation is highly sensitive to two reference
choices: the accumulation window of the ERA5 truth (a mismatched 6-hour
reference distorts the rankings) and the unit harmonization of model
precipitation that is output as a mean rate rather than a daily total;
both must be reconciled before anomaly-based skill scores are comparable. The results provide an
early, real-time assessment of AI-based S2S forecasts for the Indian
winter and establish a reproducible verification template extensible to
further systems and seasons.
\end{abstract}

\vspace{0.5em}
\noindent\textbf{Keywords:} subseasonal-to-seasonal prediction;
machine-learning weather forecasting; precipitation verification;
India; western disturbances; ensemble forecasting

\section{Introduction}\label{sec:intro}

Subseasonal-to-seasonal (S2S) prediction, covering lead times from
roughly two weeks to two months, bridges the gap between medium-range
weather forecasting and seasonal climate outlooks
\citep{vitart2017,nas2016}. Forecasts at these leads are increasingly
demanded for agricultural planning, reservoir operations, energy load
management, and disaster preparedness \citep{white2017}, yet they
remain among the most difficult products to produce skillfully: the
atmosphere has largely lost its initial-condition memory, while slowly
varying boundary forcings have not yet come to dominate
\citep{mariotti2020}. Documented sources of subseasonal predictability
include the Madden--Julian oscillation \citep[MJO;][]{zhang2005},
stratosphere--troposphere coupling, soil-moisture and snow anomalies,
and the slowly evolving state of the El Ni\~no--Southern Oscillation
(ENSO) \citep{mariotti2020,domeisen2022}. Precipitation, the variable
of greatest practical value to many users, is also the variable for
which S2S skill degrades most rapidly: hindcast-based assessments show
that anomaly correlations over most tropical and subtropical land
regions fall to modest values beyond week~2
\citep{deandrade2019,pegion2019}.

Over India, the January--March (JFM) season---boreal winter
transitioning into early spring---has received far less attention in
the S2S literature than the summer monsoon, despite its societal
importance. Winter precipitation over northern India and the Western
Himalaya is produced chiefly by western disturbances (WDs),
extratropical synoptic-scale systems embedded in the subtropical
westerly jet \citep{dimri2015,hunt2018}. WD-associated precipitation
sustains the rabi (winter) cropping season, including the wheat belt of
the Indo-Gangetic Plain, and builds the Himalayan snowpack that feeds
the Indus and Ganges river systems during the subsequent melt season
\citep{dimri2015}. Southern peninsular India can additionally receive
residual northeast-monsoon rainfall in early January, while anomalously
wet or dry winters modulate fog occurrence, cold-wave conditions, and
pre-monsoon water availability. Skillful multi-week precipitation
outlooks for JFM therefore carry substantial value for Indian
agriculture and water management, and extended-range forecasting for
the region is an operational priority of the India Meteorological
Department (IMD) and the Ministry of Earth Sciences
\citep{rao2019,pattanaik2019}.

In parallel, weather prediction is being reshaped by machine-learning
(ML) forecast systems trained on reanalysis data. Deterministic
data-driven models such as FourCastNet, Pangu-Weather, and GraphCast
now match or exceed operational numerical weather prediction (NWP) on
standard medium-range headline scores
\citep{pathak2022,bi2023,lam2023,benbouallegue2024,rasp2024};
probabilistic ML systems have followed \citep{price2025}; hybrid
approaches couple learned parameterizations to dynamical cores
\citep{kochkov2024}; and operational centers have begun running
data-driven systems alongside their NWP suites \citep{lang2024}.
Extension of these methods to the subseasonal range is more recent:
FuXi-S2S, an ML system producing 42-day ensemble forecasts, has been
reported to outperform conventional S2S models on several global
metrics \citep{chen2024}, and commercial providers now issue AI-based
S2S products \citep{spire2026}. Published evaluations of ML S2S
systems have, however, emphasized global or large-scale
diagnostics---for example, MJO prediction and midlatitude circulation
indices---while regional precipitation skill over densely populated,
hydrologically sensitive domains such as India remains largely
undocumented.

A further gap concerns the evaluation setting. Most S2S skill
estimates derive from multi-year hindcasts
\citep{deandrade2019,pegion2019}, which provide robust statistics but
may not reflect the configuration of a system as actually operated in a
given season---particularly for rapidly evolving AI models. Real-time
(or quasi-real-time) verification of a single season offers a
complementary, if statistically noisier, view: it tests the systems as
users experienced them, with operational initial conditions and
operational model versions. In this study we adopt this real-time
perspective for JFM 2026, benchmarking two AI-based systems---Spire
AI-S2S and FuXi-S2S---against two established operational dynamical
systems---the ECMWF extended-range ensemble and NCEP CFSv2---for weekly
precipitation over India.

Specifically, we address four questions:
\begin{enumerate}
  \item How does weekly skill for precipitation and for the 500-hPa
        circulation over India decay with lead time (weeks 1--6) in
        each system during JFM 2026?
  \item Is SPIRE competitive with the operational dynamical systems
        (ECMWF, NCEP), and is it superior to the other AI system
        (FuXi-S2S), and at which leads?
  \item Does skill for the large-scale circulation (Z500) translate
        into skill for precipitation?
  \item How is skill organized across the IMD homogeneous regions, and
        what systematic biases characterize each system?
\end{enumerate}
We verify ensemble-mean forecasts from 13 weekly initializations
against ERA5 using the anomaly correlation coefficient (ACC),
root-mean-square error (RMSE), and mean bias. The remainder of the
paper is organized as follows. Section~\ref{sec:data} describes the
verification data and study domain, Section~\ref{sec:method} the
processing methodology, Section~\ref{sec:models} the forecast systems,
and Section~\ref{sec:verif} the verification framework. Results are
presented in Section~\ref{sec:results} and discussed in
Section~\ref{sec:discussion}; limitations, future work, and conclusions
follow in Sections~\ref{sec:limitations}--\ref{sec:conclusions}.

\section{Data}\label{sec:data}

\subsection{Verification reference: ERA5}\label{sec:era5}

We use the fifth-generation ECMWF reanalysis
\citep[ERA5;][]{hersbach2020} as the verification reference, accessed
through the Analysis-Ready, Cloud-Optimized ERA5 archive
\citep[ARCO-ERA5;][]{carver2023}. Hourly total precipitation
(\texttt{tp}) on the native $0.25^{\circ}\times0.25^{\circ}$ grid is
summed to true 24-hour daily totals (00--00~UTC) and subsequently to the
weekly windows defined in Section~\ref{sec:weeks}; using a correctly
accumulated daily total, rather than a partial-day accumulation, is
essential for an unbiased precipitation comparison
(Section~\ref{sec:metrics}). ERA5 precipitation is a short-range forecast
product of the IFS rather than a direct observational analysis, and its
quality over the Indian region varies spatially, with the largest
uncertainties over the Western Himalaya, where gauge density is low and
orographic gradients are strong. We therefore treat ERA5 as a
dynamically consistent, temporally complete reference suitable for
anomaly-based verification, and we discuss the implications of this
choice---together with a planned cross-check against the IMD
$0.25^{\circ}$ gauge-based rainfall analysis \citep{pai2014}---in
Sections~\ref{sec:limitations} and \ref{sec:future}.
\ph{If a second high-resolution verification dataset (e.g., IMD gridded
rainfall or IMERG) is added to the final analysis, document it here and
in Table~\ref{tab:verif}.} ERA5 data through 10~May~2026 are used, which
close the final week-6 verification window of the latest (26~March)
initialization (Table~\ref{tab:verif}).

\subsection{Study domain}\label{sec:domain}

The study domain covers the Indian region, $5^{\circ}$--$38^{\circ}$N,
$65^{\circ}$--$100^{\circ}$E, with all aggregate skill metrics computed
over \emph{land points only}: ocean grid cells (Arabian Sea, Bay of
Bengal) are excluded using a coastline-based land mask, since they
otherwise dominate the area average and distort precipitation skill.
Area weighting by the cosine of latitude is applied throughout. Skill
is additionally partitioned into four India Meteorological Department
(IMD) homogeneous-region boxes---Northwest, Central, South Peninsula,
and East/Northeast India (Fig.~\ref{fig:domain})---each intersected
with the land mask.

\begin{figure}[t]
  \centering
  \includegraphics[width=0.8\linewidth]{figs/fig01_domain.pdf}
  \caption{The India study domain ($5^{\circ}$--$38^{\circ}$N,
  $65^{\circ}$--$100^{\circ}$E). Land is shaded tan and ocean blue; only
  land points are retained for the area-aggregated scores. The coloured
  boxes denote the four IMD homogeneous regions used in the regional
  analysis (Northwest, Central, South Peninsula, East/Northeast India),
  and the dashed box marks the full verification domain. All 13 weekly
  forecasts span 1~January--26~March~2026.}
  \label{fig:domain}
\end{figure}

\subsection{Forecast data}\label{sec:fcstdata}

Weekly initialized forecasts are collected from the four systems
summarized in Table~\ref{tab:models} for the 13 initialization dates
listed in Section~\ref{sec:weeks}. ECMWF and NCEP forecasts are
obtained from the WWRP/WCRP S2S database \citep{vitart2017}, using the
perturbed-forecast ensembles. Spire AI-S2S forecasts are provided by
Spire Global under a research arrangement. FuXi-S2S forecasts are
generated from the publicly released model, initialized from ERA5
analyses at 00~UTC on each start date \citep{chen2024}.

\begin{table}[t]
  \caption{Verification dataset and protocol summary.}
  \label{tab:verif}
  \centering
  \small
  \begin{tabular}{ll}
    \toprule
    Property & Specification \\
    \midrule
    Reference dataset        & ERA5 \citep{hersbach2020} \\
    Variables                & Precipitation, Z500, 2-m temperature \\
    Native resolution        & ERA5 hourly/daily (ARCO for tp, t2m) \\
    Verification grid        & $1.5^{\circ}$ common grid (conservative remapping) \\
    Temporal aggregation     & Hourly/daily $\rightarrow$ daily (00--00 UTC) $\rightarrow$ weekly \\
    Verification period      & 1 January -- 6 May 2026 \\
    Observed climatology     & In-sample 2026 study-period daily mean (Section~\ref{sec:anom}) \\
    Spatial mask             & Indian land points, $\cos\varphi$ area weights \\
    Secondary reference      & IMD $0.25^{\circ}$ gauge analysis \citep{pai2014} --- planned \\
    \bottomrule
  \end{tabular}
\end{table}

\section{Methodology}\label{sec:method}

\subsection{Initialization dates and weekly lead windows}\label{sec:weeks}

Thirteen initializations at weekly intervals span the JFM 2026 season:
1, 8, 15, 22, and 29 January; 5, 12, 19, and 26 February; and 5, 12,
19, and 26 March 2026 (all at \ph{00~UTC --- confirm per system}).
Forecast precipitation is aggregated into nonoverlapping weekly
accumulation windows: week~1 corresponds to forecast days 1--7, week~2
to days 8--14, week~3 to days 15--21, and week~4 to days 22--28, where
day~1 denotes the first 24~h after initialization. The verifying
periods collectively span 1~January--\ph{22~April}~2026. \ph{If any
system provides output at a different time resolution or
initialization time (e.g., daily-mean rates rather than accumulations,
or non-00~UTC cycles), document the alignment and conversion here.}
Where forecast horizons permit (Table~\ref{tab:models}), weeks 5--6 may
be added in a future revision.

\subsection{Common grid}\label{sec:grid}

All forecasts and the ERA5 reference are remapped to a common
$1.5^{\circ}$ latitude--longitude grid (matching the coarsest delivered
system, the S2S archive grid) prior to verification, so that no
high-resolution system is penalized for representing fine scales the
others cannot. Bilinear interpolation is used for the continuous fields;
precipitation is remapped so as to preserve area-integrated totals.

\subsection{Climatologies and anomaly definition}\label{sec:anom}

Both forecast and observed anomalies are defined by removing a single
common climatology: the ERA5 JFM-2026 seasonal mean field, computed
over the verification period and applied identically to every system
and to the ERA5 reference. Multi-decadal ERA5 climatologies and
system-specific reforecast (hindcast) climatologies were not available
for the real-time configurations evaluated here, so a lead-dependent
model-drift correction could not be applied. Two properties make this
choice defensible for the present comparison. First, the identical
baseline is subtracted from all four systems and from the reference, so
no system is advantaged relative to another. Second, our headline
metric is the spatially \emph{centered} pattern correlation
(Eq.~\ref{eq:acc}), which measures agreement in the \emph{shape} of the
anomaly field and is therefore insensitive to lead-dependent mean
offsets that an in-sample seasonal climatology leaves uncorrected. The
principal residual risk---that an in-sample seasonal mean can modestly
inflate absolute anomaly amplitudes---bears on RMSE rather than on the
pattern correlation, and is carried as a limitation
(Section~\ref{sec:limitations}).

\subsection{Ensemble treatment}\label{sec:ens}

Verification targets the ensemble mean of each system, computed
identically across all four: the SPIRE provider-supplied ensemble mean,
the FuXi-S2S 11-member mean, and the ECMWF and NCEP perturbed-member
(PF) means (100 and 15 members, respectively; member counts in
Table~\ref{tab:models}). Ensemble means from unequal ensemble sizes are
not strictly comparable, because larger ensembles damp unpredictable
variance more effectively; this asymmetry favours the larger ECMWF and
SPIRE ensembles at longer leads and is carried as a caveat
(Section~\ref{sec:limitations}). Probabilistic verification is deferred
to future work (Section~\ref{sec:future}).

\section{Forecast Models}\label{sec:models}

Table~\ref{tab:models} summarizes the four systems evaluated, as
delivered through the WWRP/WCRP S2S database and the two AI providers.

\begin{table}[t]
  \caption{Forecast systems evaluated in this study, as configured for
  the JFM~2026 forecasts used here. The ``Members'' column gives the
  ensemble actually used in verification; resolution is the native grid
  as delivered (all systems are conservatively regridded to a common
  $1.5^{\circ}$ grid for verification).}
  \label{tab:models}
  \centering
  \small
  \setlength{\tabcolsep}{4pt}
  \begin{tabular}{llllll}
    \toprule
    System & Type & Resolution & Members & Horizon & Reforecasts \\
    \midrule
    ECMWF ext-range & Dynamical, coupled & $1.5^{\circ}$ (IFS Tco319) & 100 (PF) & 46 d & on-the-fly, 20 yr \\
    NCEP CFSv2      & Dynamical, coupled & $1.5^{\circ}$ (T126)       & 15 (PF)  & 44 d & 1999--2010 \\
    Spire AI-S2S    & AI (data-driven)   & $0.5^{\circ}$             & ens.\ mean\,$+$\,spread & 46 d & --- \\
    FuXi-S2S        & AI (data-driven)   & $1.5^{\circ}$             & 11 & 42 d & --- \\
    \bottomrule
  \end{tabular}
\end{table}

\subsection{ECMWF extended-range ensemble}\label{sec:ecmwf}

The ECMWF extended-range configuration \citep{vitart2014} provides
ensemble forecasts to day~46 at a native IFS Tco319 ($\approx$36~km)
resolution, archived on the S2S database at $1.5^{\circ}$. From each
00~UTC initialization we use the 100-member perturbed-forecast ensemble
(the unperturbed control is not included in the ensemble mean).
On-the-fly reforecasts supply the lead-dependent model climatology that
would be used for provider-native anomalies; here, for a level playing
field across systems, anomalies are instead formed against a common ERA5
climatology (Section~\ref{sec:anom}).

\subsection{NCEP CFSv2}\label{sec:ncep}

The NCEP Climate Forecast System version~2 is a fully coupled
atmosphere--ocean--land system, operational since 2011
\citep{saha2014}, with a T126 ($\approx$100~km) atmosphere, archived on
the S2S database at $1.5^{\circ}$ with a 44-day horizon. From each
initialization we use the 15-member perturbed-forecast ensemble. The
1999--2010 CFSv2 reforecast set provides the provider-native model
climatology.

\subsection{Spire AI-S2S}\label{sec:spire}

Spire AI-S2S is a commercially produced, ML-based subseasonal
forecasting system \citep{spire2026}. It is delivered on a
$0.5^{\circ}$ global grid at daily resolution to a 46-day horizon. Rather
than individual ensemble members, the provider supplies ensemble summary
products---the ensemble mean and standard deviation, together with a set
of distribution percentiles and threshold-exceedance probabilities---which
we use directly: the ensemble mean for deterministic scores and the
mean--spread pair for the probabilistic verification of
Section~\ref{sec:metrics}. Forecasts used in this study cover the 13
weekly JFM~2026 initializations at 00~UTC.

\subsection{FuXi-S2S}\label{sec:fuxi}

FuXi-S2S is an ML subseasonal forecasting system based on the FuXi
architecture \citep{chen2024}, producing daily global forecasts to
day~42. While the published system supports a large perturbed ensemble,
the forecasts available to us for JFM~2026 comprise an 11-member ensemble
on a $1.5^{\circ}$ grid, initialized from ERA5 fields at 00~UTC; we use
the 11-member mean for deterministic scores and the member spread for the
probabilistic verification. The FuXi-S2S precipitation channel is
provided as a mean rate and is converted to a daily total before scoring
(Section~\ref{sec:metrics}). Because FuXi-S2S (like Spire) is initialized
directly from ERA5 whereas the dynamical systems use their own analyses,
the AI systems share the verifying reanalysis's initial state; the
implications are discussed in Section~\ref{sec:limitations}.

\section{Verification Framework}\label{sec:verif}

Let $f'_{i,t,\ell}$ denote the ensemble-mean forecast anomaly at land
grid point $i$, initialization $t$ ($t=1,\dots,T$; $T=13$), and weekly
lead $\ell\in\{1,2,3,4\}$, and let $o'_{i,t,\ell}$ denote the
corresponding ERA5 anomaly for the same verifying week. Area weights
are $w_i=\cos\varphi_i$, with sums taken over the land mask of
Fig.~\ref{fig:domain}.

\subsection{Metrics}\label{sec:metrics}

The anomaly pattern correlation (PCC; we use this term in preference to
the unqualified ``ACC'' to emphasize the spatial centering) at lead
$\ell$ is the cosine-weighted, spatially \emph{centered} Pearson
correlation of the anomaly fields, computed per initialization and
averaged over the season,
\begin{equation}
  \mathrm{PCC}(\ell) \;=\; \frac{1}{T}\sum_{t=1}^{T}
  \frac{\sum_i w_i\,(f'_{i,t,\ell}-\overline{f'}_{t,\ell})\,
                     (o'_{i,t,\ell}-\overline{o'}_{t,\ell})}
       {\sqrt{\sum_i w_i\,(f'_{i,t,\ell}-\overline{f'}_{t,\ell})^{2}}\,
        \sqrt{\sum_i w_i\,(o'_{i,t,\ell}-\overline{o'}_{t,\ell})^{2}}}\,,
  \label{eq:acc}
\end{equation}
where $\overline{x}_{t,\ell}=\big(\sum_i w_i x_{i,t,\ell}\big)/\sum_i w_i$
denotes the cosine-weighted spatial mean over land points. This is the
exact estimator implemented in the verification code; the seasonal
average is taken on the raw per-initialization correlations (not on
Fisher-$z$ transforms). The root-mean-square error is
\begin{equation}
  \mathrm{RMSE}(\ell) \;=\; \frac{1}{T}\sum_{t=1}^{T}
  \sqrt{\frac{\sum_i w_i\,\bigl(f_{i,t,\ell}-o_{i,t,\ell}\bigr)^{2}}
             {\sum_i w_i}}\,,
  \label{eq:rmse}
\end{equation}
evaluated on the weekly-mean fields in their native units
(mm\,day$^{-1}$ for precipitation, m for Z500, K for temperature), and
the mean bias is
\begin{equation}
  \mathrm{Bias}(\ell) \;=\; \frac{1}{T}\sum_{t=1}^{T}
  \frac{\sum_i w_i\,\bigl(f_{i,t,\ell}-o_{i,t,\ell}\bigr)}
       {\sum_i w_i}\,,
  \label{eq:bias}
\end{equation}
in the field's native units. \textbf{Precipitation reference.} Both
variables are verified with the full metric set (PCC, RMSE, and bias).
We emphasize one methodological point that materially affects the
precipitation results: the precipitation reference must be a true
24-hour daily accumulation. An earlier version of this analysis used an
ERA5 precipitation field archived as a single 6-hour accumulation per
day, which under-represents the daily total by roughly an order of
magnitude and, because it shares the native scale and structure of the
reanalysis-trained AI output, spuriously inflated the apparent
precipitation skill of the system trained to emulate ERA5. We therefore
construct the precipitation reference from hourly ERA5 obtained from the
Analysis-Ready, Cloud-Optimized ERA5 archive
\citep[ARCO-ERA5;][]{carver2023}, summed to true 24-hour daily totals
($\approx$0.8--1.4~mm\,day$^{-1}$ domain mean over the season),
consistent in scale with the systems that report daily-total
precipitation. All precipitation scores reported below use this
corrected reference.

\textbf{Precipitation unit harmonization.} A second, independent
precipitation issue concerns the \emph{forecast} units. The FuXi-S2S
\texttt{tp} channel is archived as a mean precipitation rate
(mm\,h$^{-1}$): its values are flat across lead day (confirming a rate
rather than an accumulation) and roughly twenty times smaller than the
daily totals of the other systems, which are already in mm\,day$^{-1}$.
We therefore multiply the FuXi-S2S precipitation field by 24 to convert
it to a daily total before any scoring; SPIRE, ECMWF, and NCEP require no
such adjustment (their domain-mean week-1 precipitation, 0.82, 0.77, and
0.84~mm\,day$^{-1}$, already matches the observed 0.82). This
harmonization is not merely cosmetic for the anomaly correlation of
Eq.~(\ref{eq:acc}): because the anomaly is formed by subtracting a
climatology computed on the true daily scale, an un-converted FuXi-S2S
field yields an almost spatially uniform negative anomaly that collapses
its week-1 precipitation PCC to an artificial 0.43; after conversion the
same field scores 0.76, a physically sensible value. Pattern
correlations of the other systems and of all non-precipitation variables
are unaffected by this step.

\textbf{Point-by-point verification.} As a distribution-level complement
to the domain-aggregated scores, we also report the joint density of
forecast and observed values pooled over all 13 initializations, all six
lead weeks (forecast days 1--42), and all land grid points
(Fig.~\ref{fig:scatter}). For each variable--system pair this yields
$\approx$26{,}000 paired values, summarized by the coefficient of
determination $R^{2}$ (the squared Pearson correlation of the pooled
pairs), the mean absolute error (MAE), and the bias; precipitation and
Z500 anomaly are shown in physical units and temperature as the absolute
field. Because the leads are pooled, this diagnostic characterizes
overall agreement rather than lead-dependent skill, which is reported
separately in Fig.~\ref{fig:acc_lead}.

For the spatial skill maps (Fig.~\ref{fig:spatial_pcc}), the local
temporal pattern correlation at each grid point is the correlation,
across the $T=13$ initializations, between forecast and observed weekly
Z500 anomalies. With so few samples these maps are necessarily noisy and
are interpreted qualitatively.

\subsection{Reference forecasts}\label{sec:baselines}

Two reference forecasts contextualize all scores: (i)~climatology
(zero anomaly), for which the ACC of Eq.~(\ref{eq:acc}) is undefined
and the anomaly RMSE equals the observed anomaly standard deviation;
and (ii)~persistence, defined as the observed anomaly of the week
immediately preceding initialization, applied at all leads.
The persistence baseline is shown in Fig.~\ref{fig:mme}.

\subsection{Uncertainty estimation}\label{sec:uncertainty}

Confidence intervals for all aggregate scores are estimated by
bootstrap resampling of the 13 initialization dates (1000 resamples;
2.5th--97.5th percentile, i.e.\ 95\%, intervals). Differences between
systems are assessed by a paired bootstrap: for each resample of the
initialization dates we recompute the difference in the seasonal-mean
score between the two systems on the \emph{same} resampled dates, and we
call a difference significant when the resulting 95\% interval excludes
zero (equivalently, a two-sided $p<0.05$). Pairing on initialization date
removes the large init-to-init variability common to both systems and is
the basis for the significance stars in Fig.~\ref{fig:spire_fuxi}. Given
$T=13$, only large differences are expected to be distinguishable, and
all statements of relative skill in Section~\ref{sec:results} should be
read with this caveat.

\subsection{Interpreting deterministic scores for smooth forecasts}
\label{sec:smoothing}

Ensemble means---and, by construction, many ML forecasts---are
smoother than individual realizations of the atmosphere. Deterministic
scores such as RMSE therefore structurally reward low-variance
forecasts (the ``double penalty'' effect;
\citealp{wilks2019,benbouallegue2024}). To make this effect visible we
additionally examine the ratio of forecast to observed spatial standard
deviation: all four systems markedly under-disperse precipitation
($\sigma_f/\sigma_o\approx0.5$--0.7), confirming that the modest
precipitation RMSE is partly a smoothing artifact, so we interpret RMSE
alongside PCC and bias rather than in isolation \citep{murphy1988}.

\subsection{Probabilistic verification}\label{sec:prob_methods}

Because the four systems carry very different ensemble information
(SPIRE provides an ensemble mean and spread; FuXi-S2S, ECMWF, and NCEP
provide 11, 100, and 15 members), we place them on a common footing by
treating each as a Gaussian forecast with the ensemble mean $\mu$ and the
ensemble spread (standard deviation) $\sigma$ at each grid point and lead
\emph{day}. The continuous ranked probability score then has the
closed form
\begin{equation}
  \mathrm{CRPS}(\mu,\sigma,y) = \sigma\!\left[\omega\,(2\Phi(\omega)-1)
  + 2\phi(\omega) - \tfrac{1}{\sqrt{\pi}}\right],\quad \omega=\tfrac{y-\mu}{\sigma},
  \label{eq:crps}
\end{equation}
with $y$ the observation and $\Phi,\phi$ the standard-normal CDF and PDF
\citep{gneiting2007}. Scores are computed per day, then averaged over the
land points and the days of each lead week. The skill score
$\mathrm{CRPSS}=1-\mathrm{CRPS}/\mathrm{CRPS}_{\mathrm{clim}}$ is taken
against a Gaussian climatological forecast built from the in-sample mean
and standard deviation. Calibration is summarized by the spread-skill
ratio $\mathrm{SSR}=\overline{\sigma}/\mathrm{RMSE}(\mu)$, which equals
one for a perfectly reliable ensemble and is below one when the ensemble
is overconfident. We also assess the reliability of two threshold events
(a heavy-rain day, precipitation $>1$~mm\,day$^{-1}$, and a cold day, T2M
below the climatological lower tercile) using the Gaussian exceedance
probability. Computing all scores at the daily level makes the SPIRE
spread (a daily ensemble standard deviation) directly comparable to the
member spreads of the other systems.

% ===================================================================
% RESULTS
% ===================================================================
\section{Results}\label{sec:results}

\subsection{Large-scale circulation (Z500) skill}\label{sec:lead_skill}

SPIRE demonstrates strong subseasonal skill for the large-scale
atmospheric circulation over India (Fig.~\ref{fig:acc_lead}c,d). At
week~1 SPIRE attains a Z500 pattern correlation (PCC) of 0.91,
comparable to ECMWF (0.94) and substantially above FuXi-S2S (0.79) and
NCEP (0.80). It remains competitive at longer leads, with PCC of 0.64
and 0.48 at weeks~2 and~3 (versus 0.62 and 0.43 for ECMWF). RMSE tells
the same story (Fig.~\ref{fig:acc_lead}d): SPIRE's week-1 Z500 RMSE of
$\approx$12.0~m is second only to ECMWF (9.3~m) and below FuXi (12.5~m)
and NCEP (19.4~m); at weeks~5--6 SPIRE attains the \emph{lowest} RMSE of
all four systems (31.9~m, versus 35.2, 37.1, and 37.4~m for ECMWF, FuXi,
and NCEP). The paired SPIRE$-$FuXi comparison
(Fig.~\ref{fig:spire_fuxi}b) shows a statistically significant week-1
advantage for SPIRE ($\Delta\mathrm{PCC}={+}0.12$, $p<0.001$). Through
week~3, SPIRE is therefore competitive with the leading operational
system (ECMWF) and outperforms both FuXi and NCEP for circulation
prediction; at weeks~4--6 the four systems become statistically
indistinguishable as all approach low pattern skill (PCC${\lesssim}0.3$).

\begin{figure}[t]
  \centering
  \includegraphics[width=\linewidth]{figs/fig02_skill_horizon.pdf}
  \caption{Weekly S2S skill horizon over Indian land points, JFM 2026,
  for SPIRE, FuXi-S2S, ECMWF, and NCEP CFSv2 (all ensemble means), as a
  function of forecast lead (weeks 1--6), averaged over the 13
  initializations: PCC (left) and RMSE (right) for
  (a,b)~precipitation (mm\,day$^{-1}$), (c,d)~Z500 (m), and
  (e,f)~2-m temperature (K). Precipitation and temperature are verified
  against true 24-hour ERA5 daily fields; the temperature RMSE is
  bias-corrected (centered) to remove the systematic cold offset caused
  by sub-daily sampling in FuXi/ECMWF/NCEP (Section~\ref{sec:t2m_skill}).
  Shading is the bootstrap 95\% CI over the 13 initializations; the
  dashed line marks PCC${=}0.5$.}
  \label{fig:acc_lead}
\end{figure}

\begin{figure}[t]
  \centering
  \includegraphics[width=\linewidth]{figs/fig03_spire_vs_fuxi.pdf}
  \caption{SPIRE versus FuXi-S2S, paired head-to-head over Indian land
  points, JFM 2026, as SPIRE minus FuXi, for $\Delta$PCC (left) and
  $\Delta$RMSE (right) of (a,b)~precipitation, (c,d)~Z500, and
  (e,f)~2-m temperature, by lead week. Because both systems ingest ERA5
  initial conditions and are verified against ERA5, this is the most
  directly comparable contrast in the study. Bars are orange where SPIRE
  is better and blue where FuXi is better; whiskers are paired-bootstrap
  95\% intervals over the 13 initializations, and an asterisk marks leads
  at which the interval excludes zero. Temperature $\Delta$RMSE uses the
  bias-corrected error (Section~\ref{sec:t2m_skill}).}
  \label{fig:spire_fuxi}
\end{figure}

\subsection{Precipitation skill}\label{sec:precip_skill}

Verified against the true 24-hour ERA5 daily totals
(Section~\ref{sec:metrics}), SPIRE is the most skillful precipitation
forecast of the four systems (Fig.~\ref{fig:acc_lead}a,b). Its week-1
pattern correlation of 0.84 exceeds ECMWF (0.80), FuXi-S2S (0.76), and
NCEP (0.66), and it remains highest through week~4 (0.71, 0.52, and 0.40
at weeks 2--4). The same ordering holds for RMSE
(Fig.~\ref{fig:acc_lead}b): SPIRE has the lowest weekly precipitation
RMSE at every lead (0.92~mm\,day$^{-1}$ at week~1, versus 1.05 for ECMWF,
1.17 for FuXi-S2S, and 1.46 for NCEP), and its domain-mean precipitation
(0.82~mm\,day$^{-1}$) is essentially unbiased against the ERA5 reference
(0.82~mm\,day$^{-1}$). The point-by-point forecast--observation density
scatter (Fig.~\ref{fig:scatter}, top row) confirms this ordering: SPIRE
has the tightest alignment with ERA5 and a bias of $+0.01$~mm\,day$^{-1}$.
The paired SPIRE$-$FuXi difference is significantly positive for
precipitation through week~3 (Fig.~\ref{fig:spire_fuxi}a,b).

ECMWF and FuXi-S2S are close competitors for second place: their week-1
correlations (0.80 and 0.76) are statistically indistinguishable, and at
week~2 FuXi-S2S (0.65) marginally exceeds ECMWF (0.63). NCEP CFSv2 is the
weakest system, with the lowest correlation and highest RMSE at short
lead. FuXi-S2S carries a modest dry bias (domain-mean
0.65~mm\,day$^{-1}$, about 20\% below the observed 0.82~mm\,day$^{-1}$),
consistent with the smoothing of extreme rainfall expected of an
MSE-trained AI model. We emphasize that this competitive FuXi-S2S ranking
emerges only after the precipitation reference and units are reconciled
(Section~\ref{sec:metrics}): the FuXi-S2S \texttt{tp} channel is archived
as a mean precipitation rate (mm\,h$^{-1}$) rather than a daily total, so
it must be multiplied by 24 before comparison. Without this harmonization
the field is roughly twentyfold too small; subtracting the
true-daily-scale climatology then leaves an almost constant negative
anomaly that artificially collapses the anomaly correlation (to 0.43 at
week~1) and inflates the RMSE. Correcting the units restores FuXi-S2S to
a physically sensible, competitive precipitation forecast while leaving
SPIRE in first place.

Although SPIRE leads, absolute precipitation skill remains modest and
decays rapidly: by week~3 every system except SPIRE has PCC below 0.45,
and all approach low skill by weeks 5--6. The faster decay of
precipitation skill relative to Z500 skill (Section~\ref{sec:lead_skill})
directly addresses our third question: useful predictability of the
large-scale circulation extends well beyond the horizon of useful
precipitation predictability over India during JFM~2026, because
precipitation is governed by local and mesoscale processes that
circulation skill alone does not capture. Tables~\ref{tab:acc} and
\ref{tab:rmse} summarize the scores, and a representative
single-initialization forecast and error map for all four systems is
shown in Fig.~\ref{fig:precip_case} (Appendix; the corresponding Z500
case is Fig.~\ref{fig:z500_case}).

\begin{table}[t]
  \centering
  \caption{Land-mean, cosine-weighted pattern correlation (PCC) over
  India, averaged over the 13 JFM 2026 initializations, for
  precipitation (top) and Z500 (bottom). Best system per lead in bold.
  Precipitation is verified against true 24-hour ERA5 daily totals.}
  \label{tab:acc}
  \begin{tabular}{lcccc}
    \toprule
    System & Week 1 & Week 2 & Week 3 & Week 4 \\
    \midrule
    \multicolumn{5}{l}{\emph{Total precipitation}}\\
    SPIRE        & \textbf{0.84} & \textbf{0.71} & \textbf{0.52} & \textbf{0.40} \\
    FuXi-S2S     & 0.76 & 0.65 & 0.40 & 0.29 \\
    ECMWF        & 0.80 & 0.63 & 0.43 & 0.29 \\
    NCEP CFSv2   & 0.66 & 0.50 & 0.39 & 0.29 \\
    \midrule
    \multicolumn{5}{l}{\emph{Z500}}\\
    SPIRE        & 0.91 & \textbf{0.64} & \textbf{0.48} & 0.30 \\
    FuXi-S2S     & 0.79 & 0.50 & 0.27 & \textbf{0.38} \\
    ECMWF        & \textbf{0.94} & 0.62 & 0.43 & 0.17 \\
    NCEP CFSv2   & 0.80 & 0.47 & 0.29 & 0.20 \\
    \bottomrule
  \end{tabular}
\end{table}

\begin{table}[t]
  \centering
  \caption{Land-mean, cosine-weighted RMSE over India, averaged over the
  13 JFM 2026 initializations, for precipitation (mm\,day$^{-1}$; top)
  and Z500 (m; bottom). Best system per lead in bold.}
  \label{tab:rmse}
  \begin{tabular}{lcccc}
    \toprule
    System & Week 1 & Week 2 & Week 3 & Week 4 \\
    \midrule
    \multicolumn{5}{l}{\emph{Precipitation (mm\,day$^{-1}$)}}\\
    SPIRE        & \textbf{0.92} & \textbf{1.30} & \textbf{1.65} & \textbf{1.84} \\
    FuXi-S2S     & 1.17 & 1.58 & 1.88 & 2.02 \\
    ECMWF        & 1.05 & 1.45 & 1.81 & 2.01 \\
    NCEP CFSv2   & 1.46 & 1.81 & 2.03 & 2.30 \\
    \midrule
    \multicolumn{5}{l}{\emph{Z500 (m)}}\\
    SPIRE        & 12.1 & 21.8 & \textbf{23.9} & \textbf{26.6} \\
    FuXi-S2S     & 12.5 & 31.1 & 47.4 & 41.2 \\
    ECMWF        & \textbf{9.3} & \textbf{18.3} & 24.1 & 30.1 \\
    NCEP CFSv2   & 19.4 & 33.4 & 35.4 & 36.5 \\
    \bottomrule
  \end{tabular}
\end{table}

\subsection{Temperature skill}\label{sec:t2m_skill}

For 2-m temperature (Fig.~\ref{fig:acc_lead}e,f), SPIRE is again the most
skillful system by a wide margin: its week-1 pattern correlation of 0.80
roughly doubles that of FuXi-S2S (0.42), ECMWF (0.40), and NCEP (0.26),
and it retains positive skill (PCC~$\approx$0.35) through week~6 while
the others decay toward zero. SPIRE is also essentially unbiased in the
domain mean ($-0.2$~K), whereas FuXi-S2S, ECMWF, and NCEP each carry a
systematic cold bias of $-3.6$ to $-4.7$~K.

This common cold bias is a temporal-sampling effect rather than a pure
forecast error: SPIRE reports a daily-mean temperature, matching the
ERA5 daily-mean reference, whereas FuXi-S2S provides an instantaneous
00~UTC field (pre-dawn over India) and the ECMWF and NCEP archives
provide 6-hour maximum and minimum temperatures whose average is
weighted toward the cold overnight window. To compare amplitude errors
fairly we therefore report the bias-corrected (centered) RMSE,
$\sqrt{\mathrm{RMSE}^2-\mathrm{Bias}^2}$, which removes the systematic
offset (Fig.~\ref{fig:acc_lead}f). Even after this correction SPIRE has
the lowest temperature error at every lead (1.1~K at week~1 versus
3.0--3.8~K for the others), and because the pattern correlation is
invariant to the bias, SPIRE's PCC advantage reflects genuinely superior
representation of temperature anomalies. The paired SPIRE$-$FuXi
comparison (Fig.~\ref{fig:spire_fuxi}) is significantly in SPIRE's favour
for temperature at all short leads.

\subsection{Reference forecasts and the multi-model ensemble}
\label{sec:mme}

Figure~\ref{fig:mme} adds two references to the four systems: a
persistence forecast (the observed weekly anomaly immediately preceding
initialization, applied at all leads) and an equal-weight multi-model
ensemble (MME; the mean of the four systems). For Z500, every system
beats persistence by a wide margin---SPIRE attains PCC of 0.91 at week~1
and remains above persistence at every lead through week~6 (e.g.\ 0.48
versus 0.00 at week~3)---confirming that the circulation skill is
genuine and not a restatement of low-frequency persistence. The MME is
the strongest single product for Z500: its PCC matches the best
individual system at every lead (0.93 at week~1, 0.67 at week~2) while
achieving the lowest RMSE of any forecast, including a week-1 RMSE of
8.1~m below ECMWF's 9.3~m. The PCC gain of the MME over the best single
system is nonetheless small ($|\Delta\mathrm{PCC}|\leq0.04$ at weeks
1--4), so with $T=13$ we frame the MME advantage as a robustness and
error-reduction benefit rather than a large skill gain.

For precipitation, every system also beats persistence (PCC 0.52 at
week~1, 0.26 at week~3): SPIRE leads at 0.84, well above the persistence
baseline at all leads, confirming that its precipitation skill is
genuine and not merely a restatement of week-to-week persistence. Here,
however, the MME does \emph{not} improve on the best single system: its
precipitation PCC (0.80 at week~1) sits just below SPIRE's 0.84, and its
RMSE (1.02~mm\,day$^{-1}$) likewise exceeds SPIRE's 0.92. Because SPIRE is
already the strongest member, equal-weight averaging with three
lower-skill systems can only regress its precipitation skill toward the
ensemble mean. Equal-weight multi-model averaging therefore helps for
Z500, where the systems are of comparable quality, but not for
precipitation, where one member is clearly best; a skill-weighted
combination would be required to
realize a benefit, which $T=13$ does not support.

A complementary diagnostic is the ratio of forecast to observed
anomaly spatial standard deviation (Section~\ref{sec:smoothing}). For
Z500 this ratio falls for SPIRE from near unity at week~1 to about 0.5
by week~4, indicating that its ensemble mean is progressively smoother
than the observed field, which partly accounts for its low RMSE; FuXi-S2S
shows the opposite, with a ratio rising above 1.4, consistent with its
larger RMSE growth, while ECMWF and NCEP remain near 0.8--1.1. RMSE
differences between systems should be read in this light.

\subsection{Spatial distribution of skill}\label{sec:spatial_skill}

Figure~\ref{fig:spatial_pcc} maps the local temporal anomaly correlation
of Z500 (computed across the 13 initializations at each grid point) for
week~2. \ph{Describe the north--south gradient evident in the maps: high
correlation over northern India and the western-disturbance corridor,
decreasing toward the peninsula, and the close similarity of the SPIRE,
ECMWF, and MME patterns. With 13 samples per grid point these maps are
indicative.}

\begin{figure}[t]
  \centering
  \includegraphics[width=\linewidth]{figs/fig06_mme_persistence.pdf}
  \caption{Weekly skill horizon including the equal-weight multi-model
  ensemble (MME, black) and the persistence baseline (grey dotted),
  for (a)~precipitation PCC, (b)~Z500 PCC, and (c)~Z500 RMSE, over
  Indian land points, JFM 2026. Shading on the MME is the bootstrap
  95\% CI over the 13 initializations.}
  \label{fig:mme}
\end{figure}

\begin{figure*}[t]
  \centering
  \includegraphics[width=\linewidth]{figs/fig07_spatial_pcc_z500.pdf}
  \caption{Local temporal anomaly correlation of Z500 at each grid
  point, computed across the 13 JFM~2026 initializations, for week~2,
  for each system and the MME. Ocean points are masked.}
  \label{fig:spatial_pcc}
\end{figure*}

\subsection{Regional skill}\label{sec:regional_skill}

Figure~\ref{fig:ranking} breaks pattern skill down by IMD homogeneous
region (Northwest, Central, South Peninsula, East/Northeast India),
averaged over weeks 1--4, and Fig.~\ref{fig:regional_horizon} (Appendix)
shows the corresponding lead dependence. For Z500, SPIRE is the most or
equal-most skillful system over Northwest (PCC 0.60, with ECMWF 0.58)
and Central India (0.52, ECMWF 0.51), and remains competitive over the
South Peninsula (0.56; ECMWF 0.60) and the East/Northeast (0.56; ECMWF
and NCEP 0.59); circulation skill is thus high and spatially uniform
across the subcontinent for SPIRE. For precipitation (against the
corrected 24-hour reference), SPIRE has the highest skill in every
region---Northwest 0.49, Central 0.59, South Peninsula 0.47, and
East/Northeast 0.67---ahead of ECMWF (0.44--0.53) and NCEP (0.30--0.51),
while FuXi-S2S is consistently the weakest (0.17--0.45). Precipitation
skill is highest over the East/Northeast and Central regions and lowest
over the Northwest, where the season's western-disturbance precipitation
is most intermittent. With four regions and 13 initializations these
estimates are indicative rather than definitive.

\subsection{Bias analysis}\label{sec:bias}

Figure~\ref{fig:bias_lead} shows the mean bias versus lead for Z500
(panel a) and precipitation (panel b, against the true 24-hour ERA5
reference). For Z500, SPIRE carries a small positive height bias
($\approx$+10~m) while FuXi-S2S and NCEP are negatively biased
($-20$ to $-35$~m), with ECMWF closest to zero (Table~\ref{tab:bias}).
For precipitation (after the FuXi-S2S unit harmonization of
Section~\ref{sec:metrics}), SPIRE, ECMWF, and NCEP are close to unbiased
in the domain mean (within $\pm$0.05~mm\,day$^{-1}$ at week~1), while
FuXi-S2S carries a modest dry bias that grows with lead (from
$-0.16$~mm\,day$^{-1}$ at week~1 to $-0.41$~mm\,day$^{-1}$ at week~4),
consistent with the tendency of MSE-trained AI models to smooth heavy
rainfall (Section~\ref{sec:precip_skill}).

\begin{figure}[t]
  \centering
  \includegraphics[width=\linewidth]{figs/fig05_bias.pdf}
  \caption{Mean bias (forecast minus ERA5) versus lead, JFM 2026, with
  bootstrap 95\% CI shading: (a)~Z500 (m) and (b)~precipitation
  (mm\,day$^{-1}$), the latter against true 24-hour ERA5 daily totals.}
  \label{fig:bias_lead}
\end{figure}

\begin{table}[t]
  \centering
  \caption{Land-mean Z500 bias (m; forecast minus ERA5), averaged over
  the 13 JFM 2026 initializations. Precipitation bias versus lead is
  shown in Fig.~\ref{fig:bias_lead}(b).}
  \label{tab:bias}
  \begin{tabular}{lcccc}
    \toprule
    System & Week 1 & Week 2 & Week 3 & Week 4 \\
    \midrule
    SPIRE        & $+10.3$ & $+11.2$ & $+7.0$ & $+6.4$ \\
    FuXi-S2S     & $-6.3$ & $-23.1$ & $-36.8$ & $-31.2$ \\
    ECMWF        & $+4.8$ & $-0.1$ & $-3.9$ & $-9.4$ \\
    NCEP CFSv2   & $-13.8$ & $-21.0$ & $-20.1$ & $-22.3$ \\
    \bottomrule
  \end{tabular}
\end{table}

\subsection{Probabilistic skill and ensemble calibration}\label{sec:prob}

The deterministic scores use only the ensemble mean; the probabilistic
verification (Section~\ref{sec:prob_methods}) asks whether each system's
\emph{uncertainty} is trustworthy. By the continuous ranked probability
skill score (Fig.~\ref{fig:crpss}), SPIRE is the most skillful
probabilistic forecast for all three variables: its week-1 CRPSS is 0.57
for precipitation, 0.47 for Z500, and 0.63 for temperature. ECMWF is a
close competitor for precipitation (0.61) but trails for the other
fields, while FuXi-S2S and NCEP fall to low or negative CRPSS by week~2,
indicating ensembles that are, probabilistically, no better than
climatology.

The cause is calibration (Fig.~\ref{fig:ssr}). The spread-skill ratio of
SPIRE is consistently the closest to one---reaching 0.7--1.1 for Z500 and
0.8--0.95 for temperature---whereas FuXi-S2S is severely overconfident,
with a week-1 spread-skill ratio of only 0.08 (precipitation), 0.14
(Z500), and 0.03 (temperature): its ensemble spread is roughly an order
of magnitude too small, so the verifying observation routinely falls
outside its envelope (visible directly in the event spread bands,
Figs.~\ref{fig:case_t2m} and \ref{fig:event_spread}). NCEP is likewise under-dispersed; ECMWF is
moderately well calibrated for Z500 but overconfident for temperature.
The reliability diagrams for a heavy-rain day and a cold day
(Fig.~\ref{fig:reliability}, Appendix) tell the same story. We note that
the very negative temperature CRPSS of the non-SPIRE systems is amplified
by the sub-daily cold bias discussed in Section~\ref{sec:t2m_skill}: an
overconfident ensemble centred several kelvin too cold is penalized
heavily by the CRPS, so part of SPIRE's temperature advantage here
reflects its daily-mean output matching the reference. For precipitation
and Z500, which are unaffected by that sampling issue, the calibration
ranking is a clean result: \emph{SPIRE is the best-calibrated system and
FuXi-S2S the most overconfident}.

\begin{figure*}[t]
  \centering
  \includegraphics[width=\linewidth]{figs/fig15_crpss.pdf}
  \caption{Continuous ranked probability skill score (CRPSS) versus lead
  for precipitation (a), Z500 (b), and 2-m temperature (c), against a
  Gaussian climatological reference, for the four systems over India,
  JFM~2026. CRPSS${>}0$ beats climatology; the dashed line marks no
  skill. Shading is the bootstrap 95\% CI over the 13 initializations.}
  \label{fig:crpss}
\end{figure*}

\begin{figure*}[t]
  \centering
  \includegraphics[width=\linewidth]{figs/fig16_spread_skill.pdf}
  \caption{Ensemble calibration: spread-skill ratio (ensemble spread
  divided by ensemble-mean RMSE) versus lead, for precipitation (a),
  Z500 (b), and 2-m temperature (c). The dashed line (ratio${=}1$) is
  perfect calibration; values below one indicate an overconfident
  (under-dispersed) ensemble. SPIRE is consistently closest to one;
  FuXi-S2S is strongly overconfident.}
  \label{fig:ssr}
\end{figure*}

\subsection{Case study: a strongly anomalous circulation week}
\label{sec:case}

To complement the aggregate statistics, Figure~\ref{fig:case} examines
the most anomalous Z500 week of the season (verifying the week of
6--12~March~2026, selected objectively as the week with the largest
root-mean-square observed Z500 anomaly over the domain), forecast at
both week-1 and week-2 lead. The observed pattern is a strong positive
height anomaly (an anomalous ridge) over northwestern India. At week-1
lead all systems reproduce the anomaly, with SPIRE, ECMWF, and FuXi-S2S
capturing its amplitude and placement closely and NCEP weaker. At
week-2 lead the systems diverge: SPIRE retains the anomaly's amplitude
and position most faithfully, ECMWF and FuXi-S2S weaken it, and NCEP
largely loses it. The case illustrates concretely the lead-dependent
skill ranking found in the aggregate scores
(Section~\ref{sec:lead_skill}).

\begin{figure*}[t]
  \centering
  \includegraphics[width=\linewidth]{figs/fig10_case_study.pdf}
  \caption{Case study of the most anomalous Z500 week of JFM~2026
  (verifying 6--12~March): observed ERA5 anomaly (left column) and each
  system's forecast at week-1 lead (top row) and week-2 lead (bottom
  row). The same verifying week is forecast at week-1 lead from its own
  initialization and at week-2 lead from the preceding initialization.
  Ocean points masked.}
  \label{fig:case}
\end{figure*}

\subsection{Week-6 case studies across variables}\label{sec:case_wk6}

To show how the systems perform at the far end of the subseasonal range,
Figures~\ref{fig:case_t2m}--\ref{fig:case_z500} present three independent
week-6 case studies, one for each variable and each on a different
initialization (chosen as the initialization whose week-6 verifying
anomaly was strongest for that field). In each figure the top row maps
the \emph{week-6 mean} anomaly (forecast days 36--42): the observed ERA5
anomaly and each system's week-6 ensemble-mean anomaly. The bottom row
shows, for each system, the All-India area-averaged anomaly through the
\emph{full} 42-day forecast---the ensemble mean and its $\pm1\sigma$
spread against the observations---with the week-6 window shaded, so that
the maps above are the average over the shaded interval. The three cases
span the range of week-6 outcomes. For the late-March warm event
(Fig.~\ref{fig:case_t2m}) SPIRE tracks the observed warming all the way
into week~6 while the other systems run cold and damp the signal---a
clear SPIRE success at six weeks. For precipitation
(Fig.~\ref{fig:case_precip}) every system retains only a weak large-scale
pattern at this range, consistent with the near-zero week-6 precipitation
skill in Table~\ref{tab:acc}. For the late-January circulation event
(Fig.~\ref{fig:case_z500}) all four systems miss the observed positive
height anomaly at week~6, illustrating that even the most predictable
field has reached its useful horizon by six weeks. Throughout, the spread
bands expose the calibration differences quantified in
Section~\ref{sec:lead_skill}: FuXi-S2S in particular shows narrow bands
that frequently exclude the verifying observation. The same three cases
broken down by IMD homogeneous region are shown in
Figures~\ref{fig:rcase_t2m}--\ref{fig:rcase_z500}, and the full SPIRE
forecast distribution (ensemble mean and percentiles) for the temperature
and precipitation cases is shown in Figures~\ref{fig:gallery_t2m} and
\ref{fig:gallery_precip}.

\begin{figure*}[t]
  \centering
  \includegraphics[width=\linewidth]{figs/fig20_casestudy_t2m.pdf}
  \caption{Week-6 case study for 2-m temperature (initialization with the
  strongest week-6 verifying T2M anomaly). Top: ERA5 anomaly and each
  system's week-6 ensemble-mean anomaly. Bottom: All-India area-averaged
  T2M anomaly versus forecast day (ERA5 in black; ensemble mean in colour
  with $\pm1\sigma$ spread); the grey band marks the week-6 window. Ocean
  masked; coastline shown.}
  \label{fig:case_t2m}
\end{figure*}

\begin{figure*}[t]
  \centering
  \includegraphics[width=\linewidth]{figs/fig21_casestudy_precip.pdf}
  \caption{As Fig.~\ref{fig:case_t2m} but for precipitation
  (mm\,day$^{-1}$), for the initialization with the strongest week-6
  verifying precipitation anomaly.}
  \label{fig:case_precip}
\end{figure*}

\begin{figure*}[t]
  \centering
  \includegraphics[width=\linewidth]{figs/fig22_casestudy_z500.pdf}
  \caption{As Fig.~\ref{fig:case_t2m} but for the Z500 anomaly (m), for
  the initialization with the strongest week-6 verifying Z500 anomaly.}
  \label{fig:case_z500}
\end{figure*}

\begin{figure*}[t]
  \centering
  \includegraphics[width=\linewidth]{figs/fig11_scatter_density.pdf}
  \caption{Point-by-point forecast--observation density scatter over India,
  pooling all 13 initializations, six lead weeks, and land grid points
  (rows: precipitation, Z500 anomaly, 2-m temperature; columns: SPIRE,
  FuXi-S2S, ECMWF, NCEP). Colour encodes point density; the solid line is
  the 1:1 reference; colour is on a logarithmic count scale. Each panel
  reports the coefficient of determination ($R^2$), mean absolute error
  (MAE), and bias. SPIRE shows the tightest alignment with ERA5 for
  precipitation and temperature and a near-zero bias throughout. FuXi-S2S
  precipitation is shown after the unit harmonization of
  Section~\ref{sec:metrics} ($\times24$, mm\,h$^{-1}\!\to$mm\,day$^{-1}$);
  it carries a modest dry tendency but is no longer collapsed. All leads
  are pooled, so the panels reflect overall agreement rather than a
  single lead time.}
  \label{fig:scatter}
\end{figure*}

\subsection{Synthesis}\label{sec:synthesis}

Taken together (Figs.~\ref{fig:acc_lead}--\ref{fig:ranking}), the
verification ranks SPIRE first overall for JFM~2026: it has the highest
precipitation and temperature skill (PCC and RMSE) at every lead with the
lowest biases, and it is statistically tied with ECMWF for the most
skillful circulation forecast while exceeding both FuXi-S2S and NCEP.
ECMWF is the strongest dynamical system and a close second; FuXi-S2S is
competitive on precipitation and circulation at short lead once its
output is correctly unit-harmonized; and NCEP is the weakest of the four
across most variables and leads. With a single season and 13 initializations these rankings are
indicative rather than definitive, but they are consistent across
variables, metrics, regions, the multi-model ensemble, and the
point-by-point forecast--observation scatter (Fig.~\ref{fig:scatter}).

\begin{figure}[t]
  \centering
  \includegraphics[width=\linewidth]{figs/fig04_regional_scorecard.pdf}
  \caption{Regional skill scorecard by IMD homogeneous region
  (Northwest, Central, South Peninsula, East/Northeast India),
  JFM 2026: mean pattern correlation over weeks 1--4 for precipitation
  (left), Z500 (centre), and 2-m temperature (right), for each system.
  Cell values are PCC; warmer colours indicate higher skill.}
  \label{fig:ranking}
\end{figure}

% ===================================================================
% DISCUSSION
% ===================================================================
\section{Discussion}\label{sec:discussion}

First, the JFM~2026 scores are broadly consistent with hindcast-based
expectations for South Asia. Multi-model S2S assessments report useful
deterministic precipitation skill typically confined to weeks~1--2, with
regime- and teleconnection-dependent windows at longer leads
\citep{deandrade2019,white2017,mariotti2020}. Our real-time results sit
within this range: precipitation pattern correlation falls below 0.5 by
week~3 for every system (Table~\ref{tab:acc}), whereas large-scale
circulation (Z500) retains useful skill several weeks longer. That the
single most skillful system here is an AI model (SPIRE) rather than ECMWF
is the principal departure from the historical picture, though with a
single season and 13~initializations this is indicative rather than
definitive.

Second, the deterministic framing structurally favors smooth forecasts.
All four systems markedly under-disperse precipitation, with
forecast-to-observed standard-deviation ratios of roughly 0.5--0.7
(Section~\ref{sec:smoothing}), so part of the low precipitation RMSE
reflects the double-penalty effect rather than superior placement of
rainfall features \citep{murphy1988,benbouallegue2024,rasp2024}. This is
precisely why we report pattern correlation and bias alongside RMSE, and
why the probabilistic diagnostics below are informative: a forecast can
win on RMSE while being poorly calibrated.

Third, the probabilistic verification adds a dimension that the
deterministic scores miss. Treating each system as a
Gaussian(ensemble-mean, ensemble-spread) forecast, SPIRE attains the
highest continuous ranked probability skill score for all three variables
and is the best calibrated, with a spread-skill ratio closest to one;
FuXi-S2S, by contrast, is severely overconfident (spread-skill ratio well
below one), so that its competitive deterministic skill does not carry
over to a competitive probabilistic skill. Calibration, not just the
ensemble mean, distinguishes these systems.

Fourth, the comparison conflates several factors: ensemble size
(100 perturbed members for ECMWF vs.\ 15 for NCEP vs.\ a provided
mean$+$spread for Spire vs.\ 11 for FuXi-S2S), initialization analyses of
differing quality, native resolution, and climatology-correction
strategy (Section~\ref{sec:anom}). In particular, both AI systems are
initialized directly from ERA5, the same reanalysis used for
verification, which may advantage them at short lead relative to the
dynamical systems that use their own analyses. Differences in scores
should therefore be attributed to ``systems as operated,'' not to
modeling paradigms in the abstract.

Fifth, the large-scale state of JFM~2026 conditions the attainable
predictability through ENSO, the MJO, and stratospheric variability,
which open and close windows of forecast opportunity
\citep{mariotti2020,domeisen2022,wheelerhendon2004,zhang2005}. A
quantitative attribution of the 2026 windows to the concurrent
large-scale state is beyond the scope of this real-time benchmark and is
left for future work with the monitoring record.

Finally, the operational context matters for India specifically: winter
precipitation governs rabi-season agriculture, Himalayan snowpack, and
water resources, and India's operational extended-range activities
provide a natural consumer for such benchmarks
\citep{rao2019,pattanaik2019,dimri2015,hunt2018}. A skillful,
well-calibrated subseasonal temperature and circulation forecast---which
SPIRE provides here through at least week~3---is directly relevant to
cold-spell and heat-stress early warning, even where precipitation
predictability remains limited.

% ===================================================================
% LIMITATIONS
% ===================================================================
\section{Limitations}\label{sec:limitations}

Several limitations qualify the conclusions. The evaluation covers a
single season with $T=13$ initializations, so sampling uncertainty is
large, score differences are often statistically indistinguishable,
and the results characterize JFM 2026 rather than climatological
skill. ERA5 is treated as truth, yet its precipitation is a model
product with documented deficiencies over complex terrain, including
the Himalaya and orographic northwest India, and verification against
gauge-based analyses can alter both magnitudes and rankings
\citep{hersbach2020,pai2014}. Only ensemble means are verified, with
unequal ensemble sizes across systems, which advantages larger
ensembles at long leads and ignores probabilistic information
entirely. The anomaly-referencing strategy is heterogeneous --- model
climatologies derive from reforecast sets of differing length and
configuration, and any in-sample correction applied where reforecasts
are unavailable optimistically biases the corrected system
(Section~\ref{sec:anom}). Operational systems may have undergone
configuration changes within or shortly before the verification
period, so the systems evaluated are snapshots. Deterministic metrics
on weekly means mask timing errors within the week and say nothing
about sharpness or reliability. Finally, AI systems trained on ERA5
and initialized from (or near) ERA5 analyses enjoy a structural
consistency advantage when verified against ERA5 itself, and
training-period overlap with the verification season cannot be fully
excluded for recently trained systems; these caveats apply with
particular force to FuXi-S2S and Spire AI-S2S
(Section~\ref{sec:models}).

% ===================================================================
% FUTURE WORK
% ===================================================================
\section{Future Work}\label{sec:future}

Natural extensions include: (i)~adding further systems as access is
secured, including FourCastNet-3 \ph{[CITE]}, the Deep Learning Earth
System Model \citep{cresswellclay2024}, and the IMD Extended Range
Forecast System \citep{pattanaik2019}, the last enabling a direct
operational comparison for India; (ii)~extending the variable set to
2-m temperature and 500-hPa geopotential height, for which weekly
predictability horizons differ markedly from precipitation;
(iii)~extending verification to weeks 5--6 where forecast ranges
permit; (iv)~probabilistic verification (CRPS, RPSS, reliability
diagrams, ROC) using full ensembles rather than means; (v)~tercile
categorical verification (above/near/below normal) with Heidke and
Gerrity scores, aligning with Indian operational convention
\citep{pattanaik2019,wilks2019}; (vi)~stratification of skill by MJO
phase \citep{wheelerhendon2004} and by the concurrent ENSO state to
identify windows of opportunity \citep{mariotti2020};
(vii)~multi-winter extension as real-time AI archives lengthen, and
companion hindcast verification where reforecasts exist;
(viii)~cross-verification against IMD gauge-based gridded rainfall
and satellite products such as IMERG to test sensitivity to the
choice of truth \citep{pai2014}; and (ix)~process-oriented
evaluation, in particular objective tracking of western disturbances
in forecasts and reanalysis \citep{hunt2018,dimri2015}.

% ===================================================================
% CONCLUSIONS
% ===================================================================
\section{Conclusions}\label{sec:conclusions}

We benchmarked two AI-based S2S systems (SPIRE AI-S2S, FuXi-S2S)
against two operational dynamical systems (ECMWF extended range, NCEP
CFSv2) for weekly precipitation and 500-hPa geopotential height over
India across 13 initializations spanning JFM 2026, verifying
ensemble-mean forecasts against ERA5 with pattern correlation, RMSE, and
bias. The main findings are: (1)~SPIRE is the most skillful
precipitation forecast at every lead (week-1 PCC 0.84 and lowest RMSE,
0.92~mm\,day$^{-1}$, with negligible domain-mean bias), ahead of ECMWF,
NCEP, and FuXi-S2S; (1b)~SPIRE is likewise the most skillful 2-m
temperature forecast (week-1 PCC 0.80, roughly double the others) and the
only system without a large cold bias; (2)~for the large-scale
circulation, SPIRE (week-1 Z500 PCC 0.91) is statistically tied with
ECMWF (0.94) and superior to FuXi-S2S and NCEP, and beats persistence at
all leads; (3)~circulation
skill substantially outlasts precipitation skill, which falls below
PCC~0.5 by week~3, so skill for the flow does not translate into
precipitation skill; (4)~an equal-weight multi-model ensemble yields the
lowest Z500 error but does not improve precipitation, where one weak
member degrades the mean; and (5)~verifying reanalysis-trained AI
precipitation against a mismatched 6-hour ERA5 reference spuriously
reversed the precipitation ranking, a methodological caution for AI-S2S
benchmarking. Within the limits of a single season, these results
support the proposition that current AI S2S systems can provide
competitive---and, for SPIRE, leading---subseasonal guidance for India,
and they establish a reproducible real-time benchmark to be extended in
subsequent seasons.

% ===================================================================
% APPENDIX
% ===================================================================
\appendix
\section{Supplementary figures}\label{app:supp}

Tables~\ref{tab:reg_rmse} and \ref{tab:reg_pcc} give the region-wise
RMSE and pattern correlation by IMD homogeneous region, averaged over
weeks 1--4 and the 13 initializations, for all three variables, and
Fig.~\ref{fig:regional_bars} presents the same information graphically.
Figure~\ref{fig:bias_maps} shows the spatial structure of the Z500 mean
anomaly bias for week~2, complementing the domain-mean values in
Section~\ref{sec:bias}. Figure~\ref{fig:regional_horizon} shows the
lead dependence of pattern correlation within each IMD homogeneous
region, complementing the weeks 1--4 scorecard of
Fig.~\ref{fig:ranking}.

\begin{table}[t]\centering
\caption{Region-wise RMSE by IMD homogeneous region (weeks 1--4 mean, 13-init average). Best system per region in bold. Temperature uses the bias-corrected (centered) RMSE (Section~\ref{sec:t2m_skill}).}\label{tab:reg_rmse}
\begin{tabular}{lcccc}\toprule
Region & SPIRE & FuXi-S2S & ECMWF & NCEP \\
\midrule
\multicolumn{5}{l}{\emph{Precipitation (mm\,day$^{-1}$)}}\\
Northwest     & \textbf{1.40} & 2.01 & 1.42 & 1.77 \\
Central       & \textbf{0.88} & 1.23 & 1.00 & 1.03 \\
S. Peninsula  & \textbf{0.71} & 0.87 & 0.75 & 0.80 \\
East/NE       & \textbf{2.08} & 3.30 & 2.49 & 2.90 \\
\midrule
\multicolumn{5}{l}{\emph{Z500 (m)}}\\
Northwest     & \textbf{22.8} & 38.4 & 23.2 & 33.8 \\
Central       & 19.4 & 24.5 & \textbf{17.5} & 29.6 \\
S. Peninsula  & 14.7 & 12.1 & \textbf{10.5} & 22.2 \\
East/NE       & 18.4 & 23.6 & \textbf{16.2} & 30.4 \\
\midrule
\multicolumn{5}{l}{\emph{T2M (K)}}\\
Northwest     & \textbf{1.64} & 3.12 & 3.33 & 4.14 \\
Central       & \textbf{0.89} & 2.06 & 1.93 & 1.97 \\
S. Peninsula  & \textbf{0.71} & 1.60 & 1.57 & 1.73 \\
East/NE       & \textbf{1.30} & 3.36 & 3.39 & 3.32 \\
\bottomrule\end{tabular}\end{table}

\begin{table}[t]\centering
\caption{Region-wise pattern correlation (PCC) by IMD homogeneous region (weeks 1--4 mean, 13-init average). Best system per region in bold.}\label{tab:reg_pcc}
\begin{tabular}{lcccc}\toprule
Region & SPIRE & FuXi-S2S & ECMWF & NCEP \\
\midrule
\multicolumn{5}{l}{\emph{Precipitation}}\\
Northwest     & \textbf{0.49} & 0.17 & 0.46 & 0.30 \\
Central       & \textbf{0.59} & 0.45 & 0.50 & 0.49 \\
S. Peninsula  & \textbf{0.47} & 0.41 & 0.44 & 0.42 \\
East/NE       & \textbf{0.67} & 0.39 & 0.53 & 0.51 \\
\midrule
\multicolumn{5}{l}{\emph{Z500}}\\
Northwest     & \textbf{0.60} & 0.39 & 0.58 & 0.44 \\
Central       & \textbf{0.52} & 0.38 & 0.51 & 0.47 \\
S. Peninsula  & 0.56 & 0.45 & \textbf{0.60} & 0.41 \\
East/NE       & 0.56 & 0.45 & 0.59 & \textbf{0.59} \\
\midrule
\multicolumn{5}{l}{\emph{T2M}}\\
Northwest     & \textbf{0.50} & 0.25 & 0.19 & 0.09 \\
Central       & \textbf{0.61} & 0.21 & 0.28 & 0.27 \\
S. Peninsula  & \textbf{0.61} & 0.30 & 0.30 & 0.30 \\
East/NE       & \textbf{0.49} & 0.13 & 0.12 & 0.07 \\
\bottomrule\end{tabular}\end{table}

\begin{figure*}[t]
  \centering
  \includegraphics[width=\linewidth]{figs/fig08_spatial_bias_z500.pdf}
  \caption{Z500 mean anomaly bias (forecast minus ERA5, m) at each grid
  point for week~2, averaged over the 13 JFM~2026 initializations, for
  each system. Ocean points are masked.}
  \label{fig:bias_maps}
\end{figure*}

\begin{figure*}[t]
  \centering
  \includegraphics[width=\linewidth]{figs/fig09_regional_horizon.pdf}
  \caption{Pattern correlation versus lead (weeks 1--6) for precipitation
  (top row) and Z500 (bottom row) within each IMD homogeneous region
  (columns), for the four systems, JFM~2026. Dotted line marks
  PCC${=}0.5$.}
  \label{fig:regional_horizon}
\end{figure*}

\begin{figure*}[t]
  \centering
  \includegraphics[width=\linewidth]{figs/fig19_regional_bars.pdf}
  \caption{Region-wise skill by IMD homogeneous region (weeks 1--4 mean,
  13-init average): pattern correlation (left column) and RMSE (right
  column) for precipitation (top), Z500 (middle), and 2-m temperature
  (bottom), with grouped bars per system. Temperature RMSE is the
  bias-corrected (centered) RMSE. This is the graphical companion to
  Tables~\ref{tab:reg_rmse} and \ref{tab:reg_pcc}.}
  \label{fig:regional_bars}
\end{figure*}

\begin{figure*}[t]
  \centering
  \includegraphics[width=\linewidth]{figs/fig12_appendix_precip_case.pdf}
  \caption{Representative single-initialization precipitation forecast and
  error over India. For the wettest week-1 of JFM~2026 by domain-mean
  precipitation (initialization 12~March~2026), the top row shows the
  ERA5 daily total and each system's week-1 forecast (mm\,day$^{-1}$;
  FuXi-S2S unit-harmonized as in Section~\ref{sec:metrics}), and the
  bottom row shows the error (forecast minus ERA5). All four systems
  capture the northeastern rainfall maximum; SPIRE has the smallest
  large-scale error. Ocean points masked; coastline shown.}
  \label{fig:precip_case}
\end{figure*}

\begin{figure*}[t]
  \centering
  \includegraphics[width=\linewidth]{figs/fig13_appendix_z500_case.pdf}
  \caption{As Fig.~\ref{fig:precip_case} but for the Z500 anomaly (m) of
  the same initialization (12~March~2026, week-1 lead): ERA5 and each
  system's forecast anomaly (top) and the error (bottom). The systems
  agree closely on the large-scale circulation, with NCEP showing the
  largest amplitude errors. Ocean points masked; coastline shown.}
  \label{fig:z500_case}
\end{figure*}

\begin{figure*}[t]
  \centering
  \includegraphics[width=\linewidth]{figs/fig23_regional_case_t2m.pdf}
  \caption{Regional breakdown of the week-6 temperature case study
  (Fig.~\ref{fig:case_t2m}): All-India and per-IMD-region area-averaged
  2-m temperature anomaly versus forecast day. ERA5 in black; the four
  systems' ensemble means in colour; SPIRE's $\pm1\sigma$ spread shaded;
  the week-6 window is grey. SPIRE tracks the observed warming most
  closely in every region.}
  \label{fig:rcase_t2m}
\end{figure*}

\begin{figure*}[t]
  \centering
  \includegraphics[width=\linewidth]{figs/fig24_regional_case_precip.pdf}
  \caption{As Fig.~\ref{fig:rcase_t2m} but for the week-6 precipitation
  case study (Fig.~\ref{fig:case_precip}).}
  \label{fig:rcase_precip}
\end{figure*}

\begin{figure*}[t]
  \centering
  \includegraphics[width=\linewidth]{figs/fig25_regional_case_z500.pdf}
  \caption{As Fig.~\ref{fig:rcase_t2m} but for the week-6 Z500 case study
  (Fig.~\ref{fig:case_z500}).}
  \label{fig:rcase_z500}
\end{figure*}

\begin{figure*}[t]
  \centering
  \includegraphics[width=\linewidth]{figs/fig26_spire_gallery_t2m.pdf}
  \caption{SPIRE forecast distribution for the week-6 temperature case
  (init 19~March~2026). Spire does not ship individual ensemble members
  but provides distribution percentiles; shown are the ERA5 anomaly, the
  SPIRE ensemble-mean anomaly, and the 1st--99th percentile anomaly maps
  (the spread analogue of a member gallery). Ocean masked; coastline
  shown.}
  \label{fig:gallery_t2m}
\end{figure*}

\begin{figure*}[t]
  \centering
  \includegraphics[width=\linewidth]{figs/fig27_spire_gallery_precip.pdf}
  \caption{As Fig.~\ref{fig:gallery_t2m} but for precipitation
  (mm\,day$^{-1}$).}
  \label{fig:gallery_precip}
\end{figure*}

\begin{figure}[t]
  \centering
  \includegraphics[width=\linewidth]{figs/fig17_reliability.pdf}
  \caption{Reliability diagrams for two probabilistic event forecasts---a
  heavy-rain day (precipitation $>1$~mm\,day$^{-1}$, left) and a cold day
  (T2M below the climatological lower tercile, right)---pooling all
  initializations, leads, and land points. Points on the diagonal are
  perfectly reliable; points below the diagonal indicate
  over-forecasting. Probabilities are the Gaussian exceedance
  probabilities of Section~\ref{sec:prob_methods}.}
  \label{fig:reliability}
\end{figure}

\begin{figure*}[t]
  \centering
  \includegraphics[width=\linewidth]{figs/fig18_event_spread.pdf}
  \caption{Ensemble spread for the most anomalous circulation event of
  JFM~2026: the All-India area-averaged Z500 anomaly (top) and
  precipitation anomaly (bottom) versus forecast day, for each system,
  with the ERA5 observation (black), the ensemble mean (colour), and the
  $\pm1\sigma$ ensemble spread (shaded). The very narrow FuXi-S2S band
  illustrates the overconfidence quantified in Fig.~\ref{fig:ssr}.}
  \label{fig:event_spread}
\end{figure*}

% ===================================================================
% BACK MATTER
% ===================================================================
\section*{Acknowledgements}

\ph{Funding sources, supervisor, institution.} This work used the
Copernicus Climate Change Service (C3S) Climate Data Store; the
results contain modified Copernicus Climate Change Service
information (2026), and neither the European Commission nor ECMWF is
responsible for any use of this information. This work is based on
S2S data; S2S is a joint initiative of the World Weather Research
Programme (WWRP) and the World Climate Research Programme (WCRP).
\ph{Acknowledge Spire data provision per agreement.} Computations
were performed on \ph{HPC system / institution}.

\section*{Data and Code Availability}

ERA5 is available from the Copernicus Climate Data Store
(\ph{DOI/URL}). ECMWF and NCEP extended-range forecasts are available
under the terms of the S2S database and the respective centers'
real-time dissemination policies \citep{vitart2017}. FuXi-S2S model
weights and inference code are available at \ph{URL}. Spire AI-S2S
forecasts were provided by Spire Global under \ph{agreement; state
access conditions}. All verification and plotting code is available
at \ph{GitHub URL / Zenodo DOI}.

% ===================================================================
% REFERENCES
% ===================================================================
% NOTE: Embedded thebibliography for arXiv self-containment.
% Entries marked %TODO must be verified against the published record
% before submission (volume/pages/DOIs may be incomplete or wrong).
\begin{thebibliography}{99}

\bibitem[{Ben~Bouall\`egue et~al.(2024)}]{benbouallegue2024}
Ben~Bouall\`egue, Z., and Coauthors, 2024: The rise of data-driven
weather forecasting: A first statistical assessment of machine
learning--based weather forecasts in an operational-like context.
\textit{Bull. Amer. Meteor. Soc.}, \textbf{105}, E864--E883. %TODO verify pages/DOI

\bibitem[{Bi et~al.(2023)}]{bi2023}
Bi, K., L. Xie, H. Zhang, X. Chen, X. Gu, and Q. Tian, 2023: Accurate
medium-range global weather forecasting with 3D neural networks.
\textit{Nature}, \textbf{619}, 533--538.

\bibitem[{Carver and Merose(2023)}]{carver2023}
Carver, R., and A. Merose, 2023: ARCO-ERA5: An Analysis-Ready,
Cloud-Optimized reanalysis dataset. \textit{Proc.\ 22nd Conf.\ on
Artificial Intelligence for Environmental Science}, Amer.\ Meteor.\
Soc. %TODO verify citation details

\bibitem[{Chen et~al.(2024)}]{chen2024}
Chen, L., X. Zhong, H. Li, J. Wu, B. Lu, D. Chen, S.-P. Xie, L. Wu,
Q. Chao, C. Lin, Z. Hu, and Y. Qi, 2024: A machine learning model
that outperforms conventional global subseasonal forecast models.
\textit{Nat. Commun.}, \textbf{15}, 6425. %TODO verify article number/author list

\bibitem[{Cresswell-Clay et~al.(2024)}]{cresswellclay2024}
Cresswell-Clay, N., and Coauthors, 2024: A deep learning Earth system
model for stable and efficient simulation of the current climate.
arXiv preprint, arXiv:2409.16247. %TODO verify identifier/journal status

\bibitem[{de~Andrade et~al.(2019)}]{deandrade2019}
de~Andrade, F.~M., C.~A.~S. Coelho, and I.~F.~A. Cavalcanti, 2019:
Global precipitation hindcast quality assessment of the Subseasonal
to Seasonal (S2S) prediction project models. \textit{Climate Dyn.},
\textbf{52}, 5451--5475.

\bibitem[{Dimri et~al.(2015)}]{dimri2015}
Dimri, A.~P., D. Niyogi, A.~P. Barros, J. Ridley, U.~C. Mohanty,
T. Yasunari, and D.~R. Sikka, 2015: Western disturbances: A review.
\textit{Rev. Geophys.}, \textbf{53}, 225--246.

\bibitem[{Domeisen et~al.(2022)}]{domeisen2022}
Domeisen, D.~I.~V., and Coauthors, 2022: Advances in the subseasonal
prediction of extreme events: Relevant case studies across the globe.
\textit{Bull. Amer. Meteor. Soc.}, \textbf{103}, E1473--E1501. %TODO verify pages

\bibitem[{Gneiting and Raftery(2007)}]{gneiting2007}
Gneiting, T., and A.~E. Raftery, 2007: Strictly proper scoring rules,
prediction, and estimation. \textit{J. Amer. Statist. Assoc.},
\textbf{102}, 359--378.

\bibitem[{Hersbach et~al.(2020)}]{hersbach2020}
Hersbach, H., and Coauthors, 2020: The ERA5 global reanalysis.
\textit{Quart. J. Roy. Meteor. Soc.}, \textbf{146}, 1999--2049.

\bibitem[{Hunt et~al.(2018)}]{hunt2018}
Hunt, K.~M.~R., A.~G. Turner, and L.~C. Shaffrey, 2018: The evolution,
seasonality and impacts of western disturbances. \textit{Quart. J.
Roy. Meteor. Soc.}, \textbf{144}, 278--290.

\bibitem[{Kochkov et~al.(2024)}]{kochkov2024}
Kochkov, D., and Coauthors, 2024: Neural general circulation models
for weather and climate. \textit{Nature}, \textbf{632}, 1060--1066.

\bibitem[{Lam et~al.(2023)}]{lam2023}
Lam, R., and Coauthors, 2023: Learning skillful medium-range global
weather forecasting. \textit{Science}, \textbf{382}, 1416--1421.

\bibitem[{Lang et~al.(2024)}]{lang2024}
Lang, S., and Coauthors, 2024: AIFS---ECMWF's data-driven forecasting
system. arXiv preprint, arXiv:2406.01465. %TODO verify journal status

\bibitem[{Mariotti et~al.(2020)}]{mariotti2020}
Mariotti, A., and Coauthors, 2020: Windows of opportunity for skillful
forecasts subseasonal to seasonal and beyond. \textit{Bull. Amer.
Meteor. Soc.}, \textbf{101}, E608--E625.

\bibitem[{Murphy(1988)}]{murphy1988}
Murphy, A.~H., 1988: Skill scores based on the mean square error and
their relationships to the correlation coefficient. \textit{Mon. Wea.
Rev.}, \textbf{116}, 2417--2424.

\bibitem[{National Academies(2016)}]{nas2016}
National Academies of Sciences, Engineering, and Medicine, 2016:
\textit{Next Generation Earth System Prediction: Strategies for
Subseasonal to Seasonal Forecasts}. The National Academies Press,
350 pp. %TODO verify page count

\bibitem[{Pai et~al.(2014)}]{pai2014}
Pai, D.~S., L. Sridhar, M. Rajeevan, O.~P. Sreejith, N.~S. Satbhai,
and B. Mukhopadhyay, 2014: Development of a new high spatial
resolution ($0.25^{\circ} \times 0.25^{\circ}$) long period
(1901--2010) daily gridded rainfall data set over India.
\textit{Mausam}, \textbf{65}, 1--18.

\bibitem[{Pathak et~al.(2022)}]{pathak2022}
Pathak, J., and Coauthors, 2022: FourCastNet: A global data-driven
high-resolution weather model using adaptive Fourier neural
operators. arXiv preprint, arXiv:2202.11214.

\bibitem[{Pattanaik et~al.(2019)}]{pattanaik2019}
Pattanaik, D.~R., \ph{coauthors}, 2019: \ph{Extended range forecast
system of IMD --- exact title}. \ph{Journal}, \ph{vol.},
\ph{pages}. %TODO complete reference

\bibitem[{Pegion et~al.(2019)}]{pegion2019}
Pegion, K., and Coauthors, 2019: The Subseasonal Experiment (SubX):
A multimodel subseasonal prediction experiment. \textit{Bull. Amer.
Meteor. Soc.}, \textbf{100}, 2043--2060.

\bibitem[{Price et~al.(2025)}]{price2025}
Price, I., and Coauthors, 2025: Probabilistic weather forecasting
with machine learning. \textit{Nature}, \textbf{637}, 84--90. %TODO verify

\bibitem[{Rao et~al.(2019)}]{rao2019}
Rao, S.~A., and Coauthors, 2019: Monsoon Mission: A targeted activity
to improve monsoon prediction across scales. \textit{Bull. Amer.
Meteor. Soc.}, \textbf{100}, 2509--2532. %TODO verify pages

\bibitem[{Rasp et~al.(2024)}]{rasp2024}
Rasp, S., and Coauthors, 2024: WeatherBench 2: A benchmark for the
next generation of data-driven global weather models. \textit{J. Adv.
Model. Earth Syst.}, \textbf{16}, e2023MS004019.

\bibitem[{Saha et~al.(2014)}]{saha2014}
Saha, S., and Coauthors, 2014: The NCEP Climate Forecast System
version 2. \textit{J. Climate}, \textbf{27}, 2185--2208.

\bibitem[{Spire Global(2026)}]{spire2026}
Spire Global, 2026: \ph{Spire AI-S2S system description / technical
documentation; cite per data agreement}. %TODO complete reference

\bibitem[{Vitart(2014)}]{vitart2014}
Vitart, F., 2014: Evolution of ECMWF sub-seasonal forecast skill
scores. \textit{Quart. J. Roy. Meteor. Soc.}, \textbf{140},
1889--1899.

\bibitem[{Vitart et~al.(2017)}]{vitart2017}
Vitart, F., and Coauthors, 2017: The Subseasonal to Seasonal (S2S)
Prediction project database. \textit{Bull. Amer. Meteor. Soc.},
\textbf{98}, 163--173.

\bibitem[{Wheeler and Hendon(2004)}]{wheelerhendon2004}
Wheeler, M.~C., and H.~H. Hendon, 2004: An all-season real-time
multivariate MJO index: Development of an index for monitoring and
prediction. \textit{Mon. Wea. Rev.}, \textbf{132}, 1917--1932.

\bibitem[{White et~al.(2017)}]{white2017}
White, C.~J., and Coauthors, 2017: Potential applications of
subseasonal-to-seasonal (S2S) predictions. \textit{Meteor. Appl.},
\textbf{24}, 315--325.

\bibitem[{Wilks(2019)}]{wilks2019}
Wilks, D.~S., 2019: \textit{Statistical Methods in the Atmospheric
Sciences}. 4th ed. Elsevier, 840 pp. %TODO verify page count

\bibitem[{Zhang(2005)}]{zhang2005}
Zhang, C., 2005: Madden--Julian Oscillation. \textit{Rev. Geophys.},
\textbf{43}, RG2003.

\end{thebibliography}

\end{document}
